\section{Discussion}\label{sec:discussion}

\subsection{The Two Systems Fail on Different CVEs}
A difference of a few points in accuracy says little about whether two systems make the same mistakes.
On this test set they do not.
Both systems get the severity right on \nBothSev\% of CVEs and both get it wrong on \nNeitherSev\%.
On the remaining \nLoRAOnlySev\% only the fine-tuned model is right, and on \nLROnlySev\% only TF-IDF is.
One or the other is therefore correct on \nEitherSev\% of CVEs, well above the \nLoRAfSev\% of the better system alone.
The pattern repeats for the CWE (either correct on \nEitherCwe\%, against \nLoRAfCwe\%) and for the exact vector (\nEitherVec\% against \nLoRAfVec\%).
No ensemble was tested here, so this is an upper bound rather than a result, but it suggests that the two systems read descriptions differently and that a combination would be worth trying.

\subsection{What a Near Miss Looks Like}
The exact-vector rate is a severe metric: one wrong component out of eight counts as a miss.
Most misses are small.
The fine-tuned model gets every component right on \nLoRAfVecZero\% of test CVEs, exactly one wrong on \nLoRAfVecOne\%, two wrong on \nLoRAfVecTwo\%, and three or more wrong on \nLoRAfVecThreePlus\%.
TF-IDF has the same shape (\nLRfVecZero\%, \nLRfVecOne\%, \nLRfVecTwo\%, and \nLRfVecThreePlus\%).
When exactly one component is wrong, it is privileges required \nLoRAfSolePR\% of the time for the fine-tuned model and \nLRfSolePR\% for TF-IDF.
Availability impact is a distant second, at \nLoRAfSoleAvail\% and \nLRfSoleAvail\%.

\begin{table}[t]
\centering
\caption{Test CVEs discussed in Section~\ref{sec:discussion}, with the NVD labels and the output of the fine-tuned model. The last column lists the vector components the model got wrong.}
\label{tab:examples}
\footnotesize
\setlength{\tabcolsep}{4pt}
\begin{tabular}{@{}lp{5.3cm}p{1.9cm}p{1.9cm}l@{}}
\toprule
CVE & Description & NVD & Predicted & Wrong \\
\midrule
\texttt{CVE-2024-0003} & A condition exists in FlashArray Purity whereby a malicious user could use a remote administrative service to create an account on the array allowing privileged access. & 7.2 High, CWE-269 & 9.8 Critical, CWE-287 & PR \\
\texttt{CVE-2024-0008} & Web sessions in the management interface in Palo Alto Networks PAN-OS software do not expire in certain situations, making it susceptible to unauthorized access. & 8.8 High, CWE-613 & 9.8 Critical, CWE-613 & UI \\
\texttt{CVE-2024-0257} & RoboDK v5.5.4 is vulnerable to heap-based buffer overflow while processing a specific project file. The resulting memory corruption may crash the application. & 3.3 Low, CWE-122 & 7.8 High, CWE-787 & C, I, A \\
\texttt{CVE-2024-10214} & Mattermost versions 9.11.X \textless{}= 9.11.1, 9.5.x \textless{}= 9.5.9 icorrectly issues two sessions when using desktop SSO - one in the browser and one in desktop with incorrect settings. & 3.5 Low, CWE-303 & 9.8 Critical, CWE-287 & PR, UI, C, I, A \\
\texttt{CVE-2024-0204} & Authentication bypass in Fortra's GoAnywhere MFT prior to 7.4.1 allows an unauthorized user to create an admin user via the administration portal. & 9.8 Critical, CWE-425 & 9.8 Critical, CWE-287 & none \\
\texttt{CVE-2024-0149} & NVIDIA GPU Display Driver for Linux contains a vulnerability which could allow an attacker unauthorized access to files. A successful exploit of this vulnerability might lead to limited information disclosure. & 3.3 Low, CWE-125 & 3.3 Low, CWE-200 & none \\
\bottomrule
\end{tabular}
\end{table}


Table~\ref{tab:examples} shows what this looks like on individual CVEs.
CVE-2024-0003, the running example, is a privileges-required miss.
The description says a ``malicious user'' can create an account through ``a remote administrative service.''
NIST read ``administrative'' as meaning the attacker already holds high privileges and scored the CVE \nExPurityTrue{}.
The model read ``malicious user'' as an outsider, set privileges required to none, and produced \nExPurityPred{}.
Every other component matched.
The same single-word ambiguity turns a \textsc{High} into a \textsc{Critical}, which is why this one component accounts for so much of the gap between partial-match and exact-match accuracy.
CVE-2024-0008 is a user-interaction miss of the same kind.
PAN-OS management sessions that do not expire are scored \nExPanOSTrue{} by NVD, with user interaction required, and the model, seeing no mention of a user, scored them \nExPanOSPred{}.
In both cases the model's reading is defensible from the text.
The information it lacks is not in the description.

\subsection{The Low Tier}
Every system is worst on \textsc{Low}, and the reason is not only that the tier is rare.
Table~\ref{tab:tiermae} gives score MAE by true tier.
On \textsc{Critical} CVEs the fine-tuned model's score is off by \nLoRAfMaeCritical{} on average, on \textsc{High} by \nLoRAfMaeHigh{}, on \textsc{Medium} by \nLoRAfMaeMedium{}, and on \textsc{Low} by \nLoRAfMaeLow{}.
Ridge regression, which never commits to a vector, does no better on them, at \nLRdMaeLow{}.
A \textsc{Low} CVE that the model gets wrong is usually not wrong by one tier.

\begin{table}[t]
\centering
\caption{Mean absolute score error by true severity tier on the 2024 test set.}
\label{tab:tiermae}
\small
\begin{tabular}{lcccc}
\toprule
System & Critical & High & Medium & Low \\
\midrule
TF-IDF ridge regression (direct) & 1.23 & 0.80 & 0.96 & 2.91 \\
TF-IDF + LR, formula & 0.59 & 0.81 & 1.04 & 3.26 \\
Phi-3.5 + LoRA, formula & 0.53 & 0.76 & 0.92 & 2.90 \\
\bottomrule
\end{tabular}
\end{table}


The examples show why.
CVE-2024-0257 describes a heap-based buffer overflow in RoboDK that ``may crash the application.''
NVD scored it \nExRoboDKTrue{}, with limited impact, because the attack needs a crafted project file and the worst outcome mentioned is a crash.
The model scored it \nExRoboDKPred{}, with high impact on all three dimensions, which is the typical score for a heap overflow in the training data.
It applied the prior for the weakness type and ignored the limiting clause.
CVE-2024-10214, a Mattermost bug that issues two sessions during desktop single sign-on, is scored \nExMattermostTrue{} by NVD and \nExMattermostPred{} by the model, which read a session-handling flaw as an authentication bypass.
Both errors run in the same direction.
When a description contains words associated with serious weaknesses, the model scores the weakness rather than the sentence.
Class-balanced weights help TF-IDF on this tier, and the equivalent for the fine-tuned model, oversampling \textsc{Low} during training, was not tried.

\subsection{Short Descriptions}
Table~\ref{tab:length} splits the test set into quartiles of description length.
Both systems do worst on the shortest quarter, under \nLenBB{} characters, where the fine-tuned model reaches \nLoRAfSevQA\% severity accuracy and TF-IDF \nLRfSevQA\%.
Accuracy rises through the middle quartiles and levels off, with the fine-tuned model at \nLoRAfSevQC\% on the third quartile and \nLoRAfSevQD\% on the longest.
The exact-vector rate follows the same curve.
The shortest descriptions are often a product name and a weakness class with nothing else, as in CVE-2024-0057 (``.NET, .NET Framework, and Visual Studio Security Feature Bypass Vulnerability''), and for those the score has to come from the prior for that vendor and weakness.
The fine-tuned model's lead over TF-IDF is similar in every quartile, so the extra capacity does not specifically rescue short inputs.

\begin{table}[t]
\centering
\caption{Severity accuracy and exact-vector rate (\%) by quartile of description length on the 2024 test set. Quartile boundaries are 194, 285, and 444 characters.}
\label{tab:length}
\small
\setlength{\tabcolsep}{4pt}
\begin{tabular}{lcccccccc}
\toprule
 & \multicolumn{2}{c}{Shortest quartile} & \multicolumn{2}{c}{Second} & \multicolumn{2}{c}{Third} & \multicolumn{2}{c}{Longest quartile} \\
\cmidrule(lr){2-3}\cmidrule(lr){4-5}\cmidrule(lr){6-7}\cmidrule(lr){8-9}
System & Sev. & Vector & Sev. & Vector & Sev. & Vector & Sev. & Vector \\
\midrule
TF-IDF + LR, formula & 59.7 & 38.8 & 67.2 & 38.4 & 71.0 & 50.7 & 70.6 & 48.7 \\
Phi-3.5 + LoRA, formula & 61.7 & 42.0 & 69.7 & 42.0 & 74.4 & 55.7 & 73.9 & 55.5 \\
\bottomrule
\end{tabular}
\end{table}


\subsection{CWE Disagreements}
Two of the examples in Table~\ref{tab:examples} have an exact vector and a wrong CWE.
For CVE-2024-0204, an authentication bypass in GoAnywhere MFT, NVD assigned CWE-425 (direct request) and the model CWE-287 (improper authentication).
For CVE-2024-0149, an NVIDIA driver flaw that allows ``unauthorized access to files'' leading to ``limited information disclosure,'' NVD assigned CWE-125 (out-of-bounds read) and the model CWE-200 (exposure of sensitive information).
In the second case the description never mentions a bounds violation.
The NVD label records what the analyst learned from the advisory, and the model's label records what the description says.
Some share of the CWE errors reported in Table~\ref{tab:main} are of this kind, and the CWE accuracy of any description-only system is bounded by it.

\subsection{Choosing a System}
The results point to a simple division.
TF-IDF logistic regression trains in minutes, needs no GPU, and gives up a few points on every field.
Formula scoring should be used with it, since it costs nothing and fixes the rare tiers, unless precision on \textsc{Critical} matters more than recall.
The fine-tuned model is better on every field, produces consistent scores and vectors, and can name CWEs outside the frequent set.
It costs \nLoraTrainHours{} hours to train and about half a second per CVE at inference, which is fast enough for the daily volume of new CVEs.
The zero-shot model should not be used for this task.
Its written severities are wrong more often than they are right, and even the vector-derived ones trail both trained systems.
Whichever system is chosen, the \textsc{Low} tier and the privileges-required component should be treated as unreliable, and a CVE scored 9.8 by either trained system deserves a second look before it jumps the queue.
